% GENERATED FILE. Edit canonical lessons, then run npm run generate:books.
% canonical-source-hash: fnv1a64:fd3b8a3cea62a80f
% canonical-lessons: RU-C100-ring, RU-C100-photograph, RU-C100-candle, RU-C100-calendar, RU-C100-rucksack

\chapter{Ring, Photograph, Candle, Calendar, Rucksack}
\label{ch:ru-ring-photograph-candle-calendar-rucksack}
\hlchaptermodality{100}

\begin{chapteropening}
\textbf{By the end of this chapter:} \emph{I can say a ring, a photograph, a candle, a calendar and a rucksack.}

Complete the last lesson of chapter 100: I can say a ring, a photograph, a candle, a calendar and a rucksack.
\end{chapteropening}

\section[\texorpdfstring{\ru{кольцо} (koltsó)}{кольцо (koltsó)}]{\ru{кольцо} (koltsó) — a ring}

\label{lesson:RU-C100-ring}



\begin{quote}
Before the new one: say the Russian for a camera, then the Russian for a purse.
\end{quote}

\subsection*{You'll want to know: \ru{кольцо}}
\textbf{\ru{кольцо}} (\emph{koltsó}) — \textquotedblleft{}a ring\textquotedblright{}.

\begin{cousinweb}
\textbf{\ru{кольцо}} is built on old \emph{\ru{коло}}, \textquotedblleft{}a wheel, a circle\textquotedblright{} (Vasmer), the root of \emph{\ru{колесо}}, \textquotedblleft{}wheel\textquotedblright{}. Vasmer relates it to Greek \emph{kýklos}, \textquotedblleft{}circle\textquotedblright{}, and Old Norse \emph{hvel}, \textquotedblleft{}wheel\textquotedblright{}: English \emph{wheel} is a cousin.
\end{cousinweb}

\begin{grammarlens}[title={What you've built}]
One more word for this chapter.
\end{grammarlens}

\subsection*{Your turn}
\begin{itemize}
\raggedright
  \item \emph{Say it:} \emph{koltsó}
  \item \emph{Say it:} \emph{koltsó}, once more
  \item \emph{Say it:} say \emph{koltsó}
  \item \emph{Recall:} say the Russian for a camera, then the Russian for a purse, then say \emph{koltsó} again
\end{itemize}

\begin{tcolorbox}[breakable,colback=teal!4,colframe=teal!35!black,title={Before you move on}]
What does \ru{кольцо} mean? (\textquotedblleft{}A ring\textquotedblright{}.) Say it once more.
\end{tcolorbox}
\section[\texorpdfstring{\ru{фотография} (fotográfiya)}{фотография (fotográfiya)}]{\ru{фотография} (fotográfiya) — a photograph}

\label{lesson:RU-C100-photograph}



\begin{quote}
Before the new one: say the Russian for a purse, then the Russian for a ring.
\end{quote}

\subsection*{You'll want to know: \ru{фотография}}
\textbf{\ru{фотография}} (\emph{fotográfiya}) — \textquotedblleft{}a photograph\textquotedblright{}.

\begin{cousinweb}
\textbf{\ru{фотография}} is an international coinage from Greek \emph{phōs}, \emph{phōtós}, \textquotedblleft{}light\textquotedblright{}, and \emph{-graphía}, \textquotedblleft{}writing\textquotedblright{}: \textquotedblleft{}writing with light\textquotedblright{}.
\end{cousinweb}

\begin{grammarlens}[title={What you've built}]
One more word for this chapter.
\end{grammarlens}

\subsection*{Your turn}
\begin{itemize}
\raggedright
  \item \emph{Say it:} \emph{fotográfiya}
  \item \emph{Say it:} \emph{fotográfiya}, once more
  \item \emph{Say it:} say \emph{fotográfiya}
  \item \emph{Recall:} say the Russian for a purse, then the Russian for a ring, then say \emph{fotográfiya} again
\end{itemize}

\begin{tcolorbox}[breakable,colback=teal!4,colframe=teal!35!black,title={Before you move on}]
What does \ru{фотография} mean? (\textquotedblleft{}A photograph\textquotedblright{}.) Say it once more.
\end{tcolorbox}
\section[\texorpdfstring{\ru{свеча} (svechá)}{свеча (svechá)}]{\ru{свеча} (svechá) — a candle}

\label{lesson:RU-C100-candle}



\begin{quote}
Before the new one: say the Russian for a ring, then the Russian for a photograph.
\end{quote}

\subsection*{You'll want to know: \ru{свеча}}
\textbf{\ru{свеча}} (\emph{svechá}) — \textquotedblleft{}a candle\textquotedblright{}.

\begin{cousinweb}
\textbf{\ru{свеча}} $\leftarrow$ Proto-Slavic \emph{*světja}, built on \emph{\ru{свет}}, \textquotedblleft{}light\textquotedblright{} (Vasmer); Vasmer compares Sanskrit \emph{śvetya-}, \textquotedblleft{}bright\textquotedblright{}.
\end{cousinweb}

\begin{grammarlens}[title={What you've built}]
One more word for this chapter.
\end{grammarlens}

\subsection*{Your turn}
\begin{itemize}
\raggedright
  \item \emph{Say it:} \emph{svechá}
  \item \emph{Say it:} \emph{svechá}, once more
  \item \emph{Say it:} say \emph{svechá}
  \item \emph{Recall:} say the Russian for a ring, then the Russian for a photograph, then say \emph{svechá} again
\end{itemize}

\begin{tcolorbox}[breakable,colback=teal!4,colframe=teal!35!black,title={Before you move on}]
What does \ru{свеча} mean? (\textquotedblleft{}A candle\textquotedblright{}.) Say it once more.
\end{tcolorbox}
\section[\texorpdfstring{\ru{календарь} (kalendár)}{календарь (kalendár)}]{\ru{календарь} (kalendár) — a calendar}

\label{lesson:RU-C100-calendar}



\begin{quote}
Before the new one: say the Russian for a photograph, then the Russian for a candle.
\end{quote}

\subsection*{You'll want to know: \ru{календарь}}
\textbf{\ru{календарь}} (\emph{kalendár}) — \textquotedblleft{}a calendar\textquotedblright{}.

\begin{cousinweb}
\textbf{\ru{календарь}} probably came through Polish \emph{kalendarz} from Latin \emph{calendārium}, \textquotedblleft{}an account book\textquotedblright{}, from \emph{calendae}, \textquotedblleft{}the first day of the month\textquotedblright{} (Vasmer).
\end{cousinweb}

\begin{grammarlens}[title={What you've built}]
One more word for this chapter.
\end{grammarlens}

\subsection*{Your turn}
\begin{itemize}
\raggedright
  \item \emph{Say it:} \emph{kalendár}
  \item \emph{Say it:} \emph{kalendár}, once more
  \item \emph{Say it:} say \emph{kalendár}
  \item \emph{Recall:} say the Russian for a photograph, then the Russian for a candle, then say \emph{kalendár} again
\end{itemize}

\begin{tcolorbox}[breakable,colback=teal!4,colframe=teal!35!black,title={Before you move on}]
What does \ru{календарь} mean? (\textquotedblleft{}A calendar\textquotedblright{}.) Say it once more.
\end{tcolorbox}
\section[\texorpdfstring{\ru{рюкзак} (ryukzák)}{рюкзак (ryukzák)}]{\ru{рюкзак} (ryukzák) — a rucksack}

\label{lesson:RU-C100-rucksack}



\begin{quote}
Before the new one: say the Russian for a candle, then the Russian for a calendar.
\end{quote}

\subsection*{You'll want to know: \ru{рюкзак}}
\textbf{\ru{рюкзак}} (\emph{ryukzák}) — \textquotedblleft{}a rucksack\textquotedblright{}.

\begin{cousinweb}
\textbf{\ru{рюкзак}} comes from German \emph{Rucksack}: \emph{Rücken}, \textquotedblleft{}back\textquotedblright{}, with \emph{Sack}, \textquotedblleft{}bag\textquotedblright{}.
\end{cousinweb}

\begin{grammarlens}[title={What you've built}]
One more word for this chapter.
\end{grammarlens}

\subsection*{Your turn}
\begin{itemize}
\raggedright
  \item \emph{Say it:} \emph{ryukzák}
  \item \emph{Say it:} \emph{ryukzák}, once more
  \item \emph{Say it:} say \emph{ryukzák}
  \item \emph{Recall:} say the Russian for a candle, then the Russian for a calendar, then say \emph{ryukzák} again
\end{itemize}

\begin{tcolorbox}[breakable,colback=teal!4,colframe=teal!35!black,title={Before you move on}]
What does \ru{рюкзак} mean? (\textquotedblleft{}A rucksack\textquotedblright{}.) Say it once more.
\end{tcolorbox}
